% Generated by Listas/make_statements.py from pratica-lecture-03.md, item 17. Do not edit.
\textbf{Source.} Gibbons (1992), \emph{Game Theory for Applied Economists}, §2.6 Problems (Section 2.1), Problem 2.1, p.~130 {[}PDF 137{]}. Reproduced verbatim. Answer in the source: none (Gibbons prints no answers).



\begin{quote}

2.1. Suppose a parent and child play the following game, first analyzed by Becker (1974). First, the child takes an action, \(A\), that produces income for the child, \(I_C(A)\), and income for the parent, \(I_P(A)\). (Think of \(I_C(A)\) as the child's income net of any costs of the action \(A\).) Second, the parent observes the incomes \(I_C\) and \(I_P\) and then chooses a bequest, \(B\), to leave to the child. The child's payoff is \(U(I_C + B)\); the parent's is \(V(I_P - B) + kU(I_C + B)\), where \(k > 0\) reflects the parent's concern for the child's well-being. Assume that: the action is a nonnegative number, \(A \geq 0\); the income functions \(I_C(A)\) and \(I_P(A)\) are strictly concave and are maximized at \(A_C > 0\) and \(A_P > 0\), respectively; the bequest \(B\) can be positive or negative; and the utility functions \(U\) and \(V\) are increasing and strictly concave. Prove the ``Rotten Kid'' Theorem: in the backwards-induction outcome, the child chooses the action that maximizes the family's aggregate income, \(I_C(A) + I_P(A)\), even though only the parent's payoff exhibits altruism.

\end{quote}

